\chapter{Objectについて }
\hypertarget{section_object}{}\label{section_object}\index{Objectについて@{Objectについて}}
\begin{DoxyItemize}
\item シーン内にて使用することが可能 (シーンがなくなると\+Objectも消えます) \item コンポーネント(\+Component)を追加できる \item 名前設定が可能で\+GUIで情報を見ることが可能\end{DoxyItemize}
\begin{DoxyItemize}
\item // 生成方法【例】(シーン内にて) 
\begin{DoxyCode}{0}
\DoxyCodeLine{　}
\DoxyCodeLine{\ \ \ \ \textcolor{keyword}{auto}\ obj\ =\ \mbox{\hyperlink{class_scene_a8eafc112b91ebde95fd4491d409b5334}{Scene::CreateObjectPtr<Object>}}();}
\DoxyCodeLine{\ \ \ \ \textcolor{keyword}{auto}\ item\ =\ \mbox{\hyperlink{class_scene_a8eafc112b91ebde95fd4491d409b5334}{Scene::CreateObjectPtr<Item>}}();\ \ \ \ \textcolor{comment}{//<\ ItemはObjectから継承したクラス}}
\DoxyCodeLine{}
\DoxyCodeLine{\ \ \ \ \textcolor{comment}{//\ 名前を変えることも可能です。(GUI上でも名前が変わります)}}
\DoxyCodeLine{\ \ \ \ item-\/>SetName(\textcolor{stringliteral}{"{}WhiteItem"{}});}
\DoxyCodeLine{}
\DoxyCodeLine{\ \ \ \ \mbox{\hyperlink{class_scene_a8eafc112b91ebde95fd4491d409b5334}{Scene::CreateObjectPtr<Player>}}()-\/>SetName(\textcolor{stringliteral}{"{}RedPlayer"{}});\ \textcolor{comment}{//<\ 変数に捉えない場合。取得方法にて取得する。}}
\DoxyCodeLine{　}

\end{DoxyCode}
\end{DoxyItemize}
\begin{DoxyItemize}
\item // 取得方法【例】(シーン内にて) 
\begin{DoxyCode}{0}
\DoxyCodeLine{　}
\DoxyCodeLine{\ \ \ \ \textcolor{comment}{//\ 存在しているPlayerのオブジェクトを取得します}}
\DoxyCodeLine{\ \ \ \ \textcolor{keyword}{auto}\ player\ =\ \mbox{\hyperlink{class_scene_a9e143dc2e860651e852c1b45b0d9b3c3}{Scene::GetObjectPtr<Player>}}();\ \ \ \textcolor{comment}{//<\ Playerクラスが一つしかない場合はこれでOK}}
\DoxyCodeLine{}
\DoxyCodeLine{\ \ \ \ \textcolor{keyword}{auto}\ red\_player\ =\ \mbox{\hyperlink{class_scene_a9e143dc2e860651e852c1b45b0d9b3c3}{Scene::GetObjectPtr<Player>}}(\textcolor{stringliteral}{"{}RedPlayer"{}});\ \ \ \textcolor{comment}{//<\ 複数の時は名前も指定しないと別のものが取れる可能性がある}}
\DoxyCodeLine{　}

\end{DoxyCode}
\end{DoxyItemize}
\begin{DoxyItemize}
\item // 削除方法【例】(シーン内にて) 
\begin{DoxyCode}{0}
\DoxyCodeLine{　}
\DoxyCodeLine{\ \ \ \ \textcolor{comment}{//\ 存在しているオブジェクトを削除します}}
\DoxyCodeLine{}
\DoxyCodeLine{\ \ \ \ \mbox{\hyperlink{class_scene_a298d059ec0aa85b0f905519a05a3b5ce}{Scene::ReleaseObject<Player>}}();\ \ \ \textcolor{comment}{//<\ Playerクラスが一つしかない場合はこれでOK}}
\DoxyCodeLine{　}
\DoxyCodeLine{\ \ \ \ \mbox{\hyperlink{class_scene_a298d059ec0aa85b0f905519a05a3b5ce}{Scene::ReleaseObject}}(\textcolor{stringliteral}{"{}RedPlayer"{}});\ \textcolor{comment}{//<\ 名前指定で削除する場合}}

\end{DoxyCode}
\end{DoxyItemize}
\begin{DoxyItemize}
\item // 継承方法【例】 
\begin{DoxyCode}{0}
\DoxyCodeLine{　}
\DoxyCodeLine{\textcolor{keyword}{class\ }Player\ :\ \textcolor{keyword}{public}\ \mbox{\hyperlink{class_object}{Object}}}
\DoxyCodeLine{\{}
\DoxyCodeLine{\textcolor{keyword}{public}:}
\DoxyCodeLine{\ \ \ \ \textcolor{keywordtype}{bool}\ \mbox{\hyperlink{class_object_a13cd3c613f3b6a01348436993b0b73c1}{Init}}()\textcolor{keyword}{\ override}}
\DoxyCodeLine{\textcolor{keyword}{\ \ \ \ }\{}
\DoxyCodeLine{\ \ \ \ \ \ \ \ \_\_super::Init();\ \ \ \ \textcolor{comment}{//<\ Object内でInitが必要なものを定義します※忘れるとASSERTが出ます}}
\DoxyCodeLine{}
\DoxyCodeLine{\ \ \ \ \ \ \ \ \textcolor{comment}{//\ 処理化処理をかく}}
\DoxyCodeLine{}
\DoxyCodeLine{\ \ \ \ \ \ \ \ \textcolor{keywordflow}{return}\ \textcolor{keyword}{true};\ \ \ \ \ \ \ \ \textcolor{comment}{//<\ 初期化が終了したら\ trueで返してください\ (※現状機能不完全)}}
\DoxyCodeLine{\ \ \ \ \ \ \ \ \ \ \ \ \ \ \ \ \ \ \ \ \ \ \ \ \ \ \ \ \textcolor{comment}{//\ \ falseの場合は、もう一度Init()が呼ばれます。(ロード完了で次へ進むなどの場合を想定しています)}}
\DoxyCodeLine{\ \ \ \ \}}
\DoxyCodeLine{}
\DoxyCodeLine{\ \ \ \ \textcolor{keywordtype}{void}\ \mbox{\hyperlink{class_object_ab5c74337ccd7504eec24590b8c963141}{Update}}()\textcolor{keyword}{\ override}}
\DoxyCodeLine{\textcolor{keyword}{\ \ \ \ }\{}
\DoxyCodeLine{\ \ \ \ \ \ \ \ \textcolor{comment}{//\ 更新処理を行う}}
\DoxyCodeLine{\ \ \ \ \}}
\DoxyCodeLine{}
\DoxyCodeLine{\ \ \ \ \textcolor{keywordtype}{void}\ \mbox{\hyperlink{class_object_ac469dbbd3cedec6a3e8f0d2172b7429e}{Draw}}()\textcolor{keyword}{\ override}}
\DoxyCodeLine{\textcolor{keyword}{\ \ \ \ }\{}
\DoxyCodeLine{\ \ \ \ \ \ \ \ \textcolor{comment}{//\ 描画処理を行う}}
\DoxyCodeLine{\ \ \ \ \}}
\DoxyCodeLine{}
\DoxyCodeLine{\ \ \ \ \textcolor{keywordtype}{void}\ \mbox{\hyperlink{class_object_a51100b72b1cfd4055b54f91ca1d77929}{Exit}}()\textcolor{keyword}{\ override}}
\DoxyCodeLine{\textcolor{keyword}{\ \ \ \ }\{}
\DoxyCodeLine{\ \ \ \ \ \ \ \ \_\_super::Exit();\ \ \ \ \textcolor{comment}{//<\ Object内でExitの終了も呼びます　※忘れるとASSERTが出ます}}
\DoxyCodeLine{\ \ \ \ \ \ \ \ \textcolor{comment}{//\ 終了処理を行う}}
\DoxyCodeLine{\ \ \ \ \}}
\DoxyCodeLine{}
\DoxyCodeLine{\ \ \ \ \textcolor{keywordtype}{void}\ \mbox{\hyperlink{class_object_a901d6889420e10a7533cd7f543426b20}{GUI}}()\textcolor{keyword}{\ override}}
\DoxyCodeLine{\textcolor{keyword}{\ \ \ \ }\{}
\DoxyCodeLine{\ \ \ \ \ \ \ \ \_\_super::GUI();\ \ \ \ \ \textcolor{comment}{//<!\ オブジェクトのGUIを使えるようにしておく}}
\DoxyCodeLine{}
\DoxyCodeLine{\ \ \ \ \ \ \ \ \textcolor{comment}{//\ PlayerGUI}}
\DoxyCodeLine{\ \ \ \ \ \ \ \ ImGui::Begin(name\_.c\_str());}
\DoxyCodeLine{\ \ \ \ \ \ \ \ \{}
\DoxyCodeLine{\ \ \ \ \ \ \ \ \ \ \ \ \textcolor{comment}{//\ GUI処理をここに追加する}}
\DoxyCodeLine{\ \ \ \ \ \ \ \ \}}
\DoxyCodeLine{\ \ \ \ \ \ \ \ ImGui::End()}
\DoxyCodeLine{\ \ \ \ \}}
\DoxyCodeLine{\};}
\DoxyCodeLine{　}

\end{DoxyCode}
 \end{DoxyItemize}
